\chapter*{Preface}
Tome~3 is about chance. It starts with puzzles that a few lines of counting settle, and paradoxes where intuition
fails, and it climbs to the theorems that make randomness predictable: the law of large numbers, the central
limit theorem, the convergence of Markov chains and of martingales, and Brownian motion, the continuous curve
that a random walk becomes when it is seen from far away. It ends with the probabilistic method, where chance
proves that objects exist without ever building them.

It assumes Tomes~0 to~2 (references such as ``Tome~0, Theorem~6.4'' point there): mostly counting, series,
integrals, Stirling's formula and the Gaussian integral from Tome~0, and the symmetric group from Tome~2.
Exercises are graded \textbf{$\star$} warm-up, \textbf{$\star\star$} for a good afternoon, \textbf{$\star\star\star$}
gems, and every exercise has a hint and a complete solution at the end of its chapter. Each chapter has a
companion Python notebook. Probability is the part of mathematics where simulation is most natural: every
theorem of this book can be watched at work on a computer, and the notebooks also follow randomness into
computer science, from hash tables and randomised algorithms to PageRank and option pricing.

The rules of the earlier tomes hold here too. Every statement is proved, or, if quoted only for culture, labelled
as a remark that says where it is proved. Every solution starts with \emph{How to find it} before the
\emph{Write-up}. ``Pause and think'' boxes ask you to guess before you read on: in probability more than anywhere
else, a wrong guess is instructive.

There is one exception to ``every statement is proved'', and it is stated openly. Measure theory, the
foundation of modern probability, is the subject of Tome~6. This tome sets up the axioms in Chapter~1 and builds
the integral in Chapter~5, proving its main theorems, but it admits two facts whose proofs are long and belong
there: that the uniform probability on $[0,1]$ (Lebesgue measure) and infinite products of probabilities exist,
and the Fubini--Tonelli theorem on iterated integrals. Each time they are used, they are labelled
\emph{Admitted (proved in Tome~6)}.

The chapters follow the plan of the series. Chapters~1--4 live in the discrete world, where sums replace
integrals; Chapters~5 and~6 build continuous probability and prove the central limit theorem; Chapters~7--11
follow random processes in time, from random walks to Brownian motion; and Chapter~12 turns probability into a
tool for combinatorics.
